\section{Theoretical Analysis}
\label{app:theory}
This section formalizes three properties of \textsc{Seed} that mirror the requirements identified in Section~\ref{sec1_intro}: the hindsight supervision is \emph{on-policy}, \emph{dense}, and \emph{self-evolving}. Specifically, we show that (i) synchronized on-policy hindsight induces an occupancy-matched adaptive target, (ii) the OPD term provides decision-specific credit even when outcome-based group advantages are uninformative, and (iii) refreshing the shared analyzer controls the staleness of the auxiliary gradient as the policy changes. These are local statements about the structure and freshness of the \textsc{Seed} update; they do not by themselves imply monotonic return improvement, which additionally requires the generated hindsight skills to be behaviorally informative.

\paragraph{Notation.}
Consider outer iteration $k$. At the beginning of the iteration, the current checkpoint is frozen as the behavior policy
$\pi_k=\pi_{\theta_k}$
and also instantiates the trajectory analyzer $\mathcal A_k$. Let $C$ denote an ordinary autoregressive context at a valid action-token position, and let $d_k(c)$ be the normalized token-context occupancy induced by $\pi_k$ and the environment. Conditional on $C=c$, the realized rollout token $Y$ is sampled as
\[
Y\sim\pi_k(\cdot\mid c).
\]
After the complete trajectory $\tau$ is observed, the synchronized analyzer produces the hindsight skill
$S=\mathcal A_k(x_\tau)$.
We overload $H(c,S)$ to denote the skill-augmented version of $c$, with the same action-token prefix as in the ordinary context. At update initialization, define the detached skill gate
\begin{equation}
 g_k(c,v,S)
 =
 \sigma\!\left(
 \beta_{\mathrm{opd}}
 \left[
 \log\pi_k(v\mid H(c,S))
 -
 \log\pi_k(v\mid c)
 \right]
 \right).
 \label{eq:theory_realized_gate}
\end{equation}
Importantly, $S$ is generated from the \emph{complete} trajectory and can therefore depend on the realized token $Y$, the subsequent rollout, environment feedback, and analyzer decoding. To preserve this dependence, we define the conditional expected gate
\begin{equation}
 w_k(c,v)
 =
 \mathbb E\!\left[
 g_k(C,Y,S)
 \mid C=c,\,Y=v
 \right],
 \label{eq:conditional_expected_gate}
\end{equation}
where the expectation is over the future trajectory, environment randomness, and any analyzer randomness. Standard softmax policies have full support, and the sigmoid lies strictly between zero and one, so $0<w_k(c,v)<1$. We state the results at $\theta=\theta_k$, where the behavior, analyzer, teacher branch, and student branch are synchronized. The adaptive-target identity below extends to an inner optimization step by replacing Eq.~\ref{eq:theory_realized_gate} with its current detached numerical value while retaining the same behavior occupancy $d_k$.

\subsection{On-Policy Hindsight Produces an Occupancy-Matched Adaptive Target}
\label{app:theory_onpolicy}
The first result characterizes the expected OPD direction and makes explicit why the source trajectories should be generated by the current policy. Rather than imitating every sampled token equally, \textsc{Seed} reweights the current policy according to how strongly its own hindsight skill supports each action.

\begin{proposition}[Occupancy-matched hindsight target]
\label{prop:onpolicy_target}
For every context $c$, define
\begin{equation}
 Z_k(c)
 =
 \sum_{v\in\mathcal V}
 \pi_k(v\mid c)w_k(c,v),
 \qquad
 r_k(v\mid c)
 =
 \frac{\pi_k(v\mid c)w_k(c,v)}{Z_k(c)}.
 \label{eq:onpolicy_target}
\end{equation}
Then $0<Z_k(c)<1$, $r_k(\cdot\mid c)$ is a probability distribution, and the expected OPD gradient at update initialization satisfies
\begin{equation}
 \left.
 \nabla_\theta\mathcal L_{\mathrm{opd},k}(\theta)
 \right|_{\theta=\theta_k}
 =
 \mathbb E_{c\sim d_k}
 \left[
 Z_k(c)
 \left.
 \nabla_\theta
 D_{\mathrm{KL}}
 \left(
 r_k(\cdot\mid c)
 \,\middle\|\,
 \pi_\theta(\cdot\mid c)
 \right)
 \right|_{\theta=\theta_k}
 \right].
 \label{eq:onpolicy_target_gradient}
\end{equation}
Thus, the OPD update is a Monte Carlo estimate of distillation toward a skill-reweighted target on the current policy's own token-context occupancy.

Moreover, let $\mu_k$ be the joint distribution of a completed trajectory and one of its valid token samples, and let $\phi_{k,\theta}(x)$ be the detached per-token OPD gradient obtained by analyzing sample $x$ with the current analyzer $\mathcal A_k$. If
$\|\phi_{k,\theta}(x)\|_2\le G$
for all relevant $x$, then replacing current on-policy samples by samples from an earlier distribution $\mu_j$ incurs
\begin{equation}
 \left\|
 \mathbb E_{x\sim\mu_k}[\phi_{k,\theta}(x)]
 -
 \mathbb E_{x\sim\mu_j}[\phi_{k,\theta}(x)]
 \right\|_2
 \le
 2G\,\operatorname{TV}(\mu_k,\mu_j).
 \label{eq:occupancy_mismatch_tv}
\end{equation}
If $\mu_k$ is absolutely continuous with respect to $\mu_j$, Pinsker's inequality further gives
\begin{equation}
 \left\|
 \mathbb E_{\mu_k}[\phi_{k,\theta}]
 -
 \mathbb E_{\mu_j}[\phi_{k,\theta}]
 \right\|_2
 \le
 G\sqrt{2D_{\mathrm{KL}}(\mu_k\|\mu_j)}.
 \label{eq:occupancy_mismatch_kl}
\end{equation}
\end{proposition}

\begin{proof}
Because $0<w_k(c,v)<1$ and $\pi_k(\cdot\mid c)$ is normalized, Eq.~\ref{eq:onpolicy_target} gives $0<Z_k(c)<1$ and $\sum_v r_k(v\mid c)=1$. In Eq.~\ref{eq:opd_objective}, both the teacher log-probability and the gate are detached, so only the ordinary student log-probability contributes to the gradient. Conditioning first on $C=c$ and $Y=v$, and then using Eq.~\ref{eq:conditional_expected_gate}, yields
\begin{align*}
 &\left.
 \nabla_\theta\mathcal L_{\mathrm{opd},k}(\theta)
 \right|_{\theta=\theta_k}
 \\
 &\quad=
 -\mathbb E_{c\sim d_k}
 \left[
 \sum_{v\in\mathcal V}
 \pi_k(v\mid c)w_k(c,v)
 \left.
 \nabla_\theta\log\pi_\theta(v\mid c)
 \right|_{\theta=\theta_k}
 \right]
 \\
 &\quad=
 -\mathbb E_{c\sim d_k}
 \left[
 Z_k(c)
 \sum_{v\in\mathcal V}
 r_k(v\mid c)
 \left.
 \nabla_\theta\log\pi_\theta(v\mid c)
 \right|_{\theta=\theta_k}
 \right],
\end{align*}
which is Eq.~\ref{eq:onpolicy_target_gradient}, since the entropy of the detached target $r_k$ has zero gradient.

For Eq.~\ref{eq:occupancy_mismatch_tv}, use the dual representation of the Euclidean norm and the variational characterization of total variation. For any unit vector $u$,
\[
 \left|
 \mathbb E_{\mu_k}[u^\top\phi_{k,\theta}]
 -
 \mathbb E_{\mu_j}[u^\top\phi_{k,\theta}]
 \right|
 \le
 2G\,\operatorname{TV}(\mu_k,\mu_j),
\]
because $|u^\top\phi_{k,\theta}(x)|\le G$. Taking the supremum over $u$ proves Eq.~\ref{eq:occupancy_mismatch_tv}. Equation~\ref{eq:occupancy_mismatch_kl} follows from
$\operatorname{TV}(\mu_k,\mu_j)\le\sqrt{D_{\mathrm{KL}}(\mu_k\|\mu_j)/2}$.
\end{proof}

Equation~\ref{eq:onpolicy_target} separates the two roles of the first \textsc{Seed} requirement. The factor $d_k(c)$ makes supervision \emph{occupancy matched}: it focuses learning on contexts, actions, and failure modes actually visited by the current policy. Within each such context, $w_k(c,v)$ makes supervision \emph{skill selective}: actions that are more strongly supported by the trajectory-specific hindsight skill receive greater relative mass. In particular,
\[
 \frac{r_k(u\mid c)}{\pi_k(u\mid c)}
 >
 \frac{r_k(v\mid c)}{\pi_k(v\mid c)}
 \quad\Longleftrightarrow\quad
 w_k(c,u)>w_k(c,v).
\]
The mismatch bounds show the complementary point: even if the same current analyzer is applied to old trajectories, a stale behavior distribution can bias the auxiliary direction in proportion to its divergence from the current trajectory-token distribution. On-policy collection sets this data-distribution mismatch to zero at the rollout stage.

The target also admits a useful value interpretation. Let
$Q_k(c,v)$
be the expected task return after choosing token $v$ at context $c$ and following the current policy thereafter. Then
\begin{equation}
 \mathbb E_{v\sim r_k(\cdot\mid c)}[Q_k(c,v)]
 -
 \mathbb E_{v\sim\pi_k(\cdot\mid c)}[Q_k(c,v)]
 =
 \frac{\operatorname{Cov}_{v\sim\pi_k(\cdot\mid c)}\!\left(Q_k(c,v),w_k(c,v)\right)}{Z_k(c)}.
 \label{eq:value_alignment_identity}
\end{equation}
Therefore, whenever hindsight support is positively correlated with current-policy action value, the reweighted target has higher local expected value than the unweighted policy. This condition makes explicit what the theorem does and does not assume: \textsc{Seed} converts hindsight support into an on-policy target, while the empirical benefit depends on the generated skills assigning greater support to better decisions.

\subsection{Dense Skill Credit Remains Informative under Sparse or Tied Rewards}
\label{app:theory_dense}
Outcome-based RL assigns the same trajectory-level advantage to every valid action token in a rollout. Consequently, it cannot distinguish a locally useful decision from a harmful one within the same trajectory, and its reward-driven signal disappears completely when all outcomes in a group are tied. The next result shows that the skill-conditioned OPD term can remain non-degenerate in exactly this regime.

\begin{proposition}[Decision-specific signal under reward ties]
\label{prop:dense_credit}
Consider a rollout group in which every trajectory has the same outcome. Then
$A^{\mathrm{rl}}_{q,n,t,\ell}=0$
for every valid token, and the clipped reward-driven term in Eq.~\ref{eq:rl_objective} has zero gradient. Fix a current-policy context $c$, parameterize the ordinary student distribution as
$p_z(\cdot\mid c)=\operatorname{softmax}(z(c))$,
and define the conditional expected OPD loss up to detached constants by
\begin{equation}
 \overline{\mathcal L}_{\mathrm{opd},k,c}(z)
 =
 -\sum_{v\in\mathcal V}
 \pi_k(v\mid c)w_k(c,v)
 \log p_z(v\mid c).
 \label{eq:conditional_opd_loss}
\end{equation}
At update initialization, where $p_z(\cdot\mid c)=\pi_k(\cdot\mid c)$,
\begin{equation}
 \frac{\partial
 \overline{\mathcal L}_{\mathrm{opd},k,c}}
 {\partial z(c,v)}
 =
 \pi_k(v\mid c)
 \left[Z_k(c)-w_k(c,v)\right].
 \label{eq:dense_credit_gradient}
\end{equation}
Furthermore, with the inverse-policy weighted norm
\[
 \|a\|^2_{\pi_k^{-1}}
 =
 \sum_{v\in\mathcal V}
 \frac{a(v)^2}{\pi_k(v\mid c)},
\]
the gradient magnitude satisfies the exact identity
\begin{equation}
 \left\|
 \nabla_{z(c)}
 \overline{\mathcal L}_{\mathrm{opd},k,c}
 \right\|^2_{\pi_k^{-1}}
 =
 \operatorname{Var}_{v\sim\pi_k(\cdot\mid c)}
 \left[w_k(c,v)\right].
 \label{eq:gate_variance_identity}
\end{equation}
Hence the OPD gradient is nonzero if and only if the expected hindsight gate is non-constant over candidate tokens with positive current-policy probability.
\end{proposition}

\begin{proof}
If all group outcomes are identical, then
$R(\tau_q^{(n)})-\mu_q=0$
for every rollout, so the group-relative advantage in Eq.~\ref{eq:rl_objective} is zero. The clipped surrogate therefore contributes no reward-derived policy gradient; the separate KL regularizer may still contribute a gradient, but it does not provide task-outcome credit.

For the OPD term, the softmax derivative gives
\[
 \frac{\partial[-\log p_z(u\mid c)]}{\partial z(c,v)}
 =
 p_z(v\mid c)-\mathbf 1\{u=v\}.
\]
Applying this identity to Eq.~\ref{eq:conditional_opd_loss} and setting $p_z=\pi_k$ yields
\begin{align*}
 \frac{\partial
 \overline{\mathcal L}_{\mathrm{opd},k,c}}
 {\partial z(c,v)}
 &=
 \left(
 \sum_u\pi_k(u\mid c)w_k(c,u)
 \right)\pi_k(v\mid c)
 -
 \pi_k(v\mid c)w_k(c,v)
 \\
 &=
 \pi_k(v\mid c)
 \left[Z_k(c)-w_k(c,v)\right],
\end{align*}
which proves Eq.~\ref{eq:dense_credit_gradient}. Substituting this gradient into the weighted norm gives
\begin{align*}
 \left\|
 \nabla_{z(c)}
 \overline{\mathcal L}_{\mathrm{opd},k,c}
 \right\|^2_{\pi_k^{-1}}
 &=
 \sum_v
 \pi_k(v\mid c)
 \left[w_k(c,v)-Z_k(c)\right]^2
 \\
 &=
 \operatorname{Var}_{v\sim\pi_k(\cdot\mid c)}
 \left[w_k(c,v)\right],
\end{align*}
because $Z_k(c)=\mathbb E_{v\sim\pi_k}[w_k(c,v)]$. The variance is zero exactly when $w_k(c,v)$ is constant $\pi_k$-almost surely.
\end{proof}

Equations~\ref{eq:dense_credit_gradient}--\ref{eq:gate_variance_identity} formalize the dense credit supplied by \textsc{Seed}. Tokens with
$w_k(c,v)>Z_k(c)$
receive a negative logit derivative and are promoted by gradient descent, whereas tokens with
$w_k(c,v)<Z_k(c)$
are relatively suppressed. Thus, although every gate value is nonnegative, softmax normalization converts variation in hindsight support into signed relative credit over candidate decisions. The variance in Eq.~\ref{eq:gate_variance_identity} quantifies the strength of this token-level signal: the more the hindsight skill discriminates among alternatives, the larger the OPD update can be. This establishes \emph{informativeness}, not automatic correctness; whether the signal favors better actions is characterized by the value-alignment covariance in Eq.~\ref{eq:value_alignment_identity}.

The tied-reward case is the sharpest contrast, but the distinction is more general. GRPO broadcasts one scalar
$A^{\mathrm{rl}}_{q,n}$
across all valid tokens in a trajectory, while \textsc{Seed} assigns a context- and token-dependent coefficient through $w_k(c,v)$. It can therefore preserve useful local behaviors in failed trajectories and attenuate locally inefficient choices in successful trajectories, even when the terminal reward alone cannot identify those decisions.

\subsection{Self-Evolving Synchronization Controls Analyzer Staleness}
\label{app:theory_selfevolving}
The final requirement concerns the source of the hindsight itself. A fixed analyzer or one-time skill dataset may be informative early in training but can become mismatched to the capabilities, strategies, and failure modes of a later policy. We quantify this mismatch by its effect on the detached OPD gate and hence on the auxiliary gradient.

For a completed token sample $x=(\tau,c,v)$ from the current distribution $\mu_k$, let
$s_r(x)=\mathcal A_r(x_\tau)$
be the skill produced by an analyzer checkpoint from iteration $r$. Holding the current policy $\pi_k$ fixed for this comparison, define
\begin{equation}
 \Delta_{k,r}(x)
 =
 \log\pi_k(v\mid H(c,s_r(x)))
 -
 \log\pi_k(v\mid c),
 \qquad
 g_{k,r}(x)
 =
 \sigma\!\left(
 \beta_{\mathrm{opd}}\Delta_{k,r}(x)
 \right).
 \label{eq:analyzer_shift_gate}
\end{equation}
The first index identifies the policy whose probabilities are re-scored, while the second identifies the analyzer that supplies the skill. For a fixed trainable parameter $\theta$, define the resulting expected auxiliary gradient on current trajectories as
\begin{equation}
 U_k(\mathcal A_r;\theta)
 =
 -\mathbb E_{x\sim\mu_k}
 \left[
 g_{k,r}(x)
 \nabla_\theta\log\pi_\theta(v\mid c)
 \right].
 \label{eq:analyzer_gradient_definition}
\end{equation}

\begin{proposition}[Analyzer-staleness bound]
\label{prop:selfevolving_staleness}
Assume
$\|\nabla_\theta\log\pi_\theta(v\mid c)\|_2\le G$
for all relevant current-policy token samples. Then the discrepancy between the OPD gradient induced by the synchronized analyzer $\mathcal A_k$ and that induced by an earlier analyzer $\mathcal A_j$ satisfies
\begin{equation}
 \begin{aligned}
 \left\|
 U_k(\mathcal A_k;\theta)
 -
 U_k(\mathcal A_j;\theta)
 \right\|_2
 &\le
 G\,
 \mathbb E_{x\sim\mu_k}
 \left[
 |g_{k,k}(x)-g_{k,j}(x)|
 \right]
 \\
 &\le
 \frac{\beta_{\mathrm{opd}}G}{4}
 \mathbb E_{x\sim\mu_k}
 \left[
 |\Delta_{k,k}(x)-\Delta_{k,j}(x)|
 \right].
 \end{aligned}
 \label{eq:analyzer_staleness_bound}
\end{equation}
In particular, using the analyzer instantiated from the current checkpoint sets this cross-iteration analyzer mismatch to zero at the rollout-and-analysis stage.
\end{proposition}

\begin{proof}
Subtract Eq.~\ref{eq:analyzer_gradient_definition} for $r=k$ and $r=j$, apply Jensen's inequality, and use the score-norm bound:
\begin{align*}
 &\left\|
 U_k(\mathcal A_k;\theta)
 -
 U_k(\mathcal A_j;\theta)
 \right\|_2
 \\
 &\quad\le
 \mathbb E_{x\sim\mu_k}
 \left[
 |g_{k,k}(x)-g_{k,j}(x)|
 \left\|
 \nabla_\theta\log\pi_\theta(v\mid c)
 \right\|_2
 \right]
 \\
 &\quad\le
 G\,
 \mathbb E_{x\sim\mu_k}
 \left[
 |g_{k,k}(x)-g_{k,j}(x)|
 \right].
\end{align*}
Because
$\sup_z\sigma'(z)=1/4$,
the map $z\mapsto\sigma(\beta_{\mathrm{opd}}z)$ is
$\beta_{\mathrm{opd}}/4$-Lipschitz. Applying this fact pointwise to Eq.~\ref{eq:analyzer_shift_gate} proves the second inequality.
\end{proof}

Proposition~\ref{prop:selfevolving_staleness} measures analyzer staleness in the quantity that directly affects learning: the change in skill-induced log-probability shift on the current policy's own trajectories. It does not require parameter distance between analyzers to be meaningful, nor does it assume that the updated analyzer is always more accurate. Instead, it states that if an old and a current analyzer produce skills with different behavioral implications, their OPD gradients can differ accordingly. The factor $\beta_{\mathrm{opd}}/4$ also exposes a trade-off: a sharper gate more strongly separates supported and unsupported tokens, but is correspondingly more sensitive to stale skill-induced shifts.

Combining the data and analyzer effects gives a compact decomposition of the two synchronization mechanisms. At a fixed current policy $\pi_k$ and trainable parameter $\theta$, let
$U_{k,\theta}(\mu,\mathcal A_r)$
denote Eq.~\ref{eq:analyzer_gradient_definition} with samples drawn from $\mu$. Under the same score-norm bound,
\begin{equation}
 \begin{aligned}
 &\left\|
 U_{k,\theta}(\mu_k,\mathcal A_k)
 -
 U_{k,\theta}(\mu_j,\mathcal A_j)
 \right\|_2
 \\
 &\qquad\le
 2G\,\operatorname{TV}(\mu_k,\mu_j)
 +
 \frac{\beta_{\mathrm{opd}}G}{4}
 \mathbb E_{x\sim\mu_j}
 \left[
 |\Delta_{k,k}(x)-\Delta_{k,j}(x)|
 \right].
 \end{aligned}
 \label{eq:joint_staleness_decomposition}
\end{equation}
The first term is \emph{experience staleness}, controlled by collecting trajectories on-policy; the second is \emph{supervision staleness}, controlled by refreshing the analyzer from the latest shared checkpoint. At each outer iteration, \textsc{Seed} uses the synchronized pair $(\mu_k,\mathcal A_k)$, so both cross-iteration mismatch terms are reset at data generation. During the subsequent inner policy optimization, the analyzer remains fixed at $\mathcal A_k$; therefore, the claim is not that analyzer lag is identically zero at every optimizer step, but that any within-update lag is discarded when trajectories and skills are regenerated from the next policy checkpoint.

Taken together, Propositions~\ref{prop:onpolicy_target},~\ref{prop:dense_credit}, and~\ref{prop:selfevolving_staleness} provide a one-to-one theoretical counterpart to the three design requirements in Section~\ref{sec1_intro}. On-policy collection aligns the auxiliary objective with the current policy's visited distribution, dense skill-conditioned re-scoring supplies token-specific credit beyond terminal outcomes, and the self-evolving shared analyzer prevents the hindsight target from remaining permanently tied to an earlier policy checkpoint.
